\documentclass[aspectratio=169, xcolor=dvipsnames]{beamer}

% --- Theme Setup ---
\usetheme{Madrid}
\usecolortheme{default}

% --- Packages ---
\usepackage{graphicx}
\usepackage{xcolor}
\usepackage{booktabs}
\usepackage{hyperref}
\usepackage{adjustbox}
\usepackage{amssymb}
\usepackage{amsmath}
\usepackage{listings}
\usepackage{caption}

\lstset{
  basicstyle=\ttfamily\small,
  keywordstyle=\color{blue}\bfseries,
  stringstyle=\color{red},
  showstringspaces=false,
  breaklines=true
}

% --- Title Information ---
\title{Module 10}
\subtitle{Introduction to SQL/3}
\author{Milav Dabgar}
\institute{www.milav.in}
\date{\today}

\begin{document}

% Title Slide
\begin{frame}
    \titlepage
\end{frame}

\begin{frame}{Module Recap}
    \begin{itemize}[<+->]
        \item Completed the understanding of basic query structure
    \end{itemize}
\end{frame}

\begin{frame}{Module Objectives}
    \begin{block}{Goals}
        \begin{itemize}[<+->]
            \item To familiarize with set operations, null values and aggregation
        \end{itemize}
    \end{block}
\end{frame}

\begin{frame}{Module Outline}
    \begin{itemize}[<+->]
        \item Set Operations: \texttt{union}, \texttt{intersect}, \texttt{except}
        \item Null Values
        \item Aggregate Functions: \texttt{avg}, \texttt{min}, \texttt{max}, \texttt{sum}, and \texttt{count}
        \item \texttt{Group By}
        \item \texttt{Having}
        \item Null Values
    \end{itemize}
\end{frame}

\section{Set Operations}

\begin{frame}
  \vfill
  \centering
  \begin{beamercolorbox}[sep=8pt,center,shadow=true,rounded=true]{title}
    \usebeamerfont{title}Set Operations\par%
  \end{beamercolorbox}
  \vfill
\end{frame}

\begin{frame}[fragile]{Set Operations}
    \begin{itemize}
        \item Find courses that ran in Fall 2009 or in Spring 2010
        \begin{lstlisting}[language=SQL]
(select course_id from section where sem = 'Fall' and year = 2009)
union
(select course_id from section where sem = 'Spring' and year = 2010);
        \end{lstlisting}
        \item Find courses that ran in Fall 2009 and in Spring 2010
        \begin{lstlisting}[language=SQL]
(select course_id from section where sem = 'Fall' and year = 2009)
intersect
(select course_id from section where sem = 'Spring' and year = 2010);
        \end{lstlisting}
        \item Find courses that ran in Fall 2009 but not in Spring 2010
        \begin{lstlisting}[language=SQL]
(select course_id from section where sem = 'Fall' and year = 2009)
except
(select course_id from section where sem = 'Spring' and year = 2010);
        \end{lstlisting}
    \end{itemize}
\end{frame}

\begin{frame}[fragile]{Set Operations (2)}
    \begin{itemize}
        \item Find the salaries of all instructors that are less than the largest salary
        \begin{lstlisting}[language=SQL]
select distinct T.salary 
from instructor as T, instructor as S 
where T.salary < S.salary;
        \end{lstlisting}
        \item Find all the salaries of all instructors
        \begin{lstlisting}[language=SQL]
select distinct salary 
from instructor;
        \end{lstlisting}
        \item Find the largest salary of all instructors
        \begin{lstlisting}[language=SQL]
(select 'second query')
except
(select 'first query');
        \end{lstlisting}
    \end{itemize}
\end{frame}

\begin{frame}{Set Operations (3)}
    \begin{block}{Duplicates}
        \begin{itemize}[<+->]
            \item Set operations \texttt{union}, \texttt{intersect}, and \texttt{except}
            \item Each of the above operations automatically eliminates duplicates
            \item To retain all duplicates use the corresponding multiset versions \texttt{union all}, \texttt{intersect all}, and \texttt{except all}.
            \item Suppose a tuple occurs $m$ times in $r$ and $n$ times in $s$, then, it occurs:
            \begin{itemize}
                \item[$\triangleright$] \textbf{$m + n$} times in \texttt{$r$ union all $s$}
                \item[$\triangleright$] \textbf{$\min(m, n)$} times in \texttt{$r$ intersect all $s$}
                \item[$\triangleright$] \textbf{$\max(0, m - n)$} times in \texttt{$r$ except all $s$}
            \end{itemize}
        \end{itemize}
    \end{block}
\end{frame}

\section{Null Values}

\begin{frame}
  \vfill
  \centering
  \begin{beamercolorbox}[sep=8pt,center,shadow=true,rounded=true]{title}
    \usebeamerfont{title}Null Values\par%
  \end{beamercolorbox}
  \vfill
\end{frame}

\begin{frame}[fragile]{Null Values}
    \begin{itemize}
        \item It is possible for tuples to have a null value, denoted by \texttt{null}, for some of their attributes
        \item \texttt{null} signifies an unknown value or that a value does not exist
        \item The result of any arithmetic expression involving \texttt{null} is \texttt{null}
        \begin{itemize}
            \item Example: \texttt{5 + null} returns \texttt{null}
        \end{itemize}
        \item The predicate \texttt{is null} can be used to check for null values
        \begin{itemize}
            \item Example: Find all instructors whose salary is null
            \begin{lstlisting}[language=SQL]
select name 
from instructor 
where salary is null;
            \end{lstlisting}
        \end{itemize}
        \item It is not possible to test for null values with comparison operators, such as \texttt{=}, \texttt{<}, or \texttt{<>}. We need to use the \texttt{is null} and \texttt{is not null} operators instead
    \end{itemize}
\end{frame}

\begin{frame}{Null Values (2): Three Valued Logic}
    \begin{block}{Overview}
        \begin{itemize}[<+->]
            \item Three values: \textbf{true}, \textbf{false}, \textbf{unknown}
            \item Any comparison with null returns \textbf{unknown}
            \begin{itemize}
                \item[$\triangleright$] Example: $\texttt{5 < null}$ \quad or \quad $\texttt{null <> null}$ \quad or \quad $\texttt{null = null}$
            \end{itemize}
            \item Three-valued logic using the value \textbf{unknown}:
            \begin{itemize}
                \item[$\triangleright$] \textbf{OR}: \texttt{(unknown or true) = true}, \texttt{(unknown or false) = unknown}, \texttt{(unknown or unknown) = unknown}
                \item[$\triangleright$] \textbf{AND}: \texttt{(true and unknown) = unknown}, \texttt{(false and unknown) = false}, \texttt{(unknown and unknown) = unknown}
                \item[$\triangleright$] \textbf{NOT}: \texttt{(not unknown) = unknown}
            \end{itemize}
            \item `\texttt{P is unknown}' evaluates to \textbf{true} if predicate $P$ evaluates to \textbf{unknown}
            \item Result of \texttt{where} clause predicate is treated as \textbf{false} if it evaluates to \textbf{unknown}
        \end{itemize}
    \end{block}
\end{frame}

\section{Aggregate Functions}

\begin{frame}
  \vfill
  \centering
  \begin{beamercolorbox}[sep=8pt,center,shadow=true,rounded=true]{title}
    \usebeamerfont{title}Aggregate Functions\par%
  \end{beamercolorbox}
  \vfill
\end{frame}

\begin{frame}{Aggregate Functions}
    \begin{exampleblock}{Common Functions}
        \begin{itemize}[<+->]
            \item These functions operate on the multiset of values of a column of a relation, and return a value
            \begin{itemize}
                \item[$\triangleright$] \texttt{avg}: average value
                \item[$\triangleright$] \texttt{min}: minimum value
                \item[$\triangleright$] \texttt{max}: maximum value
                \item[$\triangleright$] \texttt{sum}: sum of values
                \item[$\triangleright$] \texttt{count}: number of values
            \end{itemize}
        \end{itemize}
    \end{exampleblock}
\end{frame}

\begin{frame}[fragile]{Aggregate Functions (2)}
    \begin{itemize}
        \item Find the average salary of instructors in the Computer Science department
        \begin{lstlisting}[language=SQL]
select avg(salary) 
from instructor 
where dept_name = 'Comp. Sci';
        \end{lstlisting}
        \item Find the total number of instructors who teach a course in the Spring 2010 semester
        \begin{lstlisting}[language=SQL]
select count(distinct ID) 
from teaches 
where semester = 'Spring' and year = 2010;
        \end{lstlisting}
        \item Find the number of tuples in the course relation
        \begin{lstlisting}[language=SQL]
select count(*) 
from course;
        \end{lstlisting}
    \end{itemize}
\end{frame}

\begin{frame}[fragile]{Aggregate Functions (3): Group By}
    \begin{itemize}
        \item Find the average salary of instructors in each department
        \begin{lstlisting}[language=SQL]
select dept_name, avg(salary) as avg_salary 
from instructor 
group by dept_name;
        \end{lstlisting}
    \end{itemize}

    \vspace{1em}
    \begin{minipage}{0.65\textwidth}
        \small
        \textbf{instructor} relation
        \begin{center}
            \begin{tabular}{|l|l|l|l|}
            \hline
            \textbf{ID} & \textbf{name} & \textbf{dept\_name} & \textbf{salary} \\ \hline
            76766 & Crick & Biology & 72000 \\ \hline
            45565 & Katz & Comp. Sci. & 75000 \\ \hline
            10101 & Srinivasan & Comp. Sci. & 65000 \\ \hline
            76543 & Singh & Finance & 80000 \\ \hline
            \dots & \dots & \dots & \dots \\ \hline
            33456 & Gold & Physics & 87000 \\ \hline
            22222 & Einstein & Physics & 95000 \\ \hline
            \end{tabular}
        \end{center}
    \end{minipage}%
    \begin{minipage}{0.35\textwidth}
        \small
        \textbf{Result}
        \begin{center}
            \begin{tabular}{|l|l|}
            \hline
            \textbf{dept\_name} & \textbf{avg\_salary} \\ \hline
            Biology & 72000 \\ \hline
            Comp. Sci. & 77333 \\ \hline
            Finance & 85000 \\ \hline
            \dots & \dots \\ \hline
            Physics & 91000 \\ \hline
            \end{tabular}
        \end{center}
    \end{minipage}
\end{frame}

\begin{frame}[fragile]{Aggregate Functions (4): Group By}
    \begin{itemize}
        \item Attributes in \texttt{select} clause outside of aggregate functions must appear in \texttt{group by} list
        \begin{lstlisting}[language=SQL]
/* erroneous query */ 
select dept_name, ID, avg(salary) 
from instructor 
group by dept_name;
        \end{lstlisting}
    \end{itemize}
\end{frame}

\begin{frame}[fragile]{Aggregate Functions (5): Having Clause}
    \begin{itemize}
        \item Find the names and average salaries of all departments whose average salary is greater than 42000
        \begin{lstlisting}[language=SQL]
select dept_name, avg(salary) 
from instructor 
group by dept_name 
having avg(salary) > 42000;
        \end{lstlisting}
        \item \textbf{Note}: Predicates in the \texttt{having} clause are applied \textbf{after} the formation of groups, whereas predicates in the \texttt{where} clause are applied \textbf{before} forming groups.
    \end{itemize}
\end{frame}

\begin{frame}[fragile]{Null Values and Aggregates}
    \begin{itemize}[<+->]
        \item Total all salaries
        \begin{lstlisting}[language=SQL]
select sum(salary) 
from instructor;
        \end{lstlisting}
        \item Above statement ignores null amounts
        \item Result is null if there is no non-null amount
        \item All aggregate operations except \texttt{count(*)} ignore tuples with null values on the aggregated attributes
        \item What if collection has only null values?
        \begin{itemize}
            \item[$\triangleright$] \texttt{count} returns \texttt{0}
            \item[$\triangleright$] all other aggregates return \texttt{null}
        \end{itemize}
    \end{itemize}
\end{frame}

\begin{frame}{Module Summary}
    \begin{itemize}[<+->]
        \item Completed the understanding of set operations, null values, and aggregation
    \end{itemize}
    \vspace{2em}
    \begin{center}
    \small \textit{Slides used in this presentation are borrowed from \url{http://db-book.com/} with kind permission of the authors.} \\
    \small \textit{Edited and new slides are marked with 'PPD'.}
    \end{center}
\end{frame}

\end{document}
